\printedpage{93}
% head: A. Grothendieck : Dérivateurs, Ch. XII, Édition des Universités de Montpellier II et Paris VII
% foot: Version 7 décembre 2002

I note that $\tsX$ being cofibered over $\Psi=\Phi^\circ$, $\tsX^\circ$ is fibered over $\Psi^\circ=\Phi$. More precisely, we will have
\[
\tsX^\circ = \int
\begin{tikzcd}[row sep=small, column sep=small]
 & (X_1)^\circ \\
(X_0)^\circ \arrow[ur, "\alpha_1^\circ"] \arrow[dr, "\alpha_2^\circ"'] & \\
 & (X_2)^\circ,
\end{tikzcd}
\]
\msp{page 137}
and it is thus a matter of comparing the two integrals
\[
\sX = \int
\begin{tikzcd}[row sep=small, column sep=small]
 & X_1 \\
X_0 \arrow[ur, "\alpha_1"] \arrow[dr, "\alpha_2"'] & \\
 & X_2
\end{tikzcd}
, \qquad
\tsX^\circ = \int
\begin{tikzcd}[row sep=small, column sep=small]
 & (X_1)^\circ \\
(X_0)^\circ \arrow[ur, "\alpha_1^\circ"] \arrow[dr, "\alpha_2^\circ"'] & \\
 & (X_2)^\circ
\end{tikzcd}
,
\]
in particular as objects of $\HOT_W(\Phi)$. For this, I take up again the construction of the functor
\[
X \longmapsto \Phi(X) : \Cat \longrightarrow \Cat
\]
from I, p. 72-75, but I prefer here to denote this functor by $\Sigma(X)$ (for $\Phi$ already denotes a small category $\begin{tikzcd}[row sep=tiny, column sep=small]
 & 1 \arrow[dl] \\
0 & \\
 & 2 \arrow[ul]
\end{tikzcd}$). We will use the homomorphisms functorial in $X$:
\[
\begin{tikzcd}[column sep=small]
 & \Sigma(X) \arrow[dl, "p_X"'] \arrow[dr, "q_X"] & \\
X & & X^\circ,
\end{tikzcd}
\]
which are $W$-aspheric, hence $\in W$. Moreover, $\Sigma(X^\circ)=\Sigma(X)$, $p_{X^\circ}=q_X$, $q_{X^\circ}=p_X$. Thus, if
\[
\begin{aligned}
X_* \quad &= \qquad
\begin{tikzcd}[ampersand replacement=\&, row sep=small, column sep=small]
 \& X_1 \\
X_0 \arrow[ur, "\alpha_1"] \arrow[dr, "\alpha_2"'] \& \\
 \& X_2
\end{tikzcd}
\qquad \in \quad \Hom(\Phi^\circ,\Cat)\\
X_*^\circ \quad &= \qquad
\begin{tikzcd}[ampersand replacement=\&, row sep=small, column sep=small]
 \& (X_1)^\circ \\
(X_0)^\circ \arrow[ur, "\alpha_1^\circ"] \arrow[dr, "\alpha_2^\circ"'] \& \\
 \& (X_2)^\circ
\end{tikzcd}
\qquad \in \quad \Hom(\Phi^\circ,\Cat),
\end{aligned}
\]
\msp{page 138}
we find
\[
\Sigma(X_*) =
\left(
\begin{tikzcd}[row sep=small, column sep=small]
 & \Sigma(X_1) \\
\Sigma(X_0) \arrow[ur, "\Sigma(\alpha_1)"] \arrow[dr, "\Sigma(\alpha_2)"'] & \\
 & \Sigma(X_2)
\end{tikzcd}
\right)
\]
and the arrows
\printedpage{94}
% head: A. Grothendieck : Dérivateurs, Ch. XII, Édition des Universités de Montpellier II et Paris VII
% foot: Version 7 décembre 2002
\[
\begin{tikzcd}[column sep=small]
 & \Sigma(X_*) \arrow[dl, "p_{X_*}"'] \arrow[dr, "q_{X_*}"] & \\
X_* & & X_*^\circ,
\end{tikzcd}
\]
in $\underline{\Hom}(\Phi^\circ,\Cat)$. These arrows are in $W$ term by term, so the homomorphisms in $\Cat$
\[
\begin{tikzcd}[column sep=small]
 & \int\Sigma(X_*) \arrow[dl, "\int p_{X_*}"'] \arrow[dr, "\int q_{X_*}"] & \\
\sX=\int X_* & & \int X_*^\circ=\tsX^\circ
\end{tikzcd}
\]
are in $W$.

Finally, the innocent commutative square
\[
Q \quad = \qquad
\begin{tikzcd}[row sep=small, column sep=small]
 & X_0 \arrow[dl] \arrow[dr] & \\
X_1 \arrow[dr] & & X_2 \arrow[dl] \\
 & Y &
\end{tikzcd}
\tag*{$(*)$}
\]
(and the corresponding diagram
\[
\widehat{Q} \quad = \qquad
\begin{tikzcd}[row sep=small, column sep=small]
 & X_0 \arrow[dl] \arrow[dr] & \\
X_1 \arrow[dr] & & X_2 \arrow[dl] \\
 & \sX=\int X_* \arrow[d, "\varphi_Q"] & \\
 & Y &
\end{tikzcd}
\qquad\qquad )
\tag*{$(**)$}
\]
gives rise, by transforming it
\msp{page 139}
by $X\longmapsto X^\circ$ and by $X\longmapsto\Sigma(X)$, and by taking into account the homomorphisms functorial in $X$ $\begin{tikzcd}[row sep=small, column sep=small]
 & \Sigma(X) \arrow[dl] \arrow[dr] & \\
X & & X^\circ
\end{tikzcd}$, to three diagrams of type $(**)$, i.e. modeled on the opposite category of
\[
\Gamma \quad = \qquad
\begin{tikzcd}[row sep=small, column sep=small]
 & 0 & \\
1 \arrow[ur, "\alpha_1"] & & 2 \arrow[ul, "\alpha_2"'] \\
 & 3 \arrow[ul, "i_1"] \arrow[ur, "i_2"'] & \\
 & 4 \arrow[u, "\varphi"'] &
\end{tikzcd}
\qquad
\text{\parbox[c]{0.35\textwidth}{\footnotesize (i.e. $\boldsymbol{\Delta}^1\times\boldsymbol{\Delta}^1$ (\footnotemark[88]), to which we have added an initial object)}} % [sic] auquelle -> à laquelle, grammatical form (catégorie is feminine)
\]
\footnotetext[88]{false, for the square is not commutative}
\printedpage{95}
% head: A. Grothendieck : Dérivateurs, Ch. XII, Édition des Universités de Montpellier II et Paris VII
% foot: Version 7 décembre 2002
(\footnote[89]{We have $\alpha_1i_1\neq\alpha_2i_2$, but the relation $\alpha_1i_1\varphi=\alpha_2i_2\varphi$.}), namely
\[
\widehat{Q}, \qquad \widehat{\Sigma(Q)}, \qquad \widehat{Q^\circ},
\]
and two homomorphisms of diagrams (\texttt{[in]} $\underline{\Hom}(\Gamma^\circ,\Cat)$)
\[
\begin{tikzcd}[column sep=small]
 & \widehat{\Sigma(Q)} \arrow[dl] \arrow[dr] & \\
\widehat{Q} & & \widehat{Q^\circ}
\end{tikzcd}
\qquad
\text{(deduced from }
\begin{tikzcd}[column sep=small]
 & \Sigma(Q) \arrow[dl] \arrow[dr] & \\
Q & & Q^\circ
\end{tikzcd}
\text{),}
\tag*{$(***)$}
\]
which are in $W$ term by term (i.e. are in $W_\Gamma\subseteq\Fl(\underline{\Hom}(\Gamma^\circ,\Cat))$). Explicitly, at the vertices $i=0,1,2$, resp. $4$ of $\Gamma$, $(***)$ is none other than
\[
D_i =
\left(
\begin{tikzcd}[column sep=small]
 & \Sigma(X_i) \arrow[dl, "p_{X_i}"'] \arrow[dr, "q_{X_i}"] & \\
X_i & & X_i^\circ
\end{tikzcd}
\right)
\qquad \text{resp.} \qquad
\begin{tikzcd}[column sep=small]
 & \Sigma(Y) \arrow[dl, "p_Y"'] \arrow[dr, "q_Y"] & \\
Y & & Y^\circ,
\end{tikzcd}
\]
at vertex 3 it is
\[
\begin{tikzcd}[column sep=small]
 & \int\Sigma(X_*) \arrow[dl, "\int p_{X_*}"'] \arrow[dr, "\int q_{X_*}"] & \\
\sX=\int X_* & & \int X_*^\circ=\tsX^\circ,
\end{tikzcd}
\]
obtained by integration of the arrows
\[
\begin{tikzcd}[column sep=small]
 & \Sigma(X_*) \arrow[dl, "p_{X_*}"'] \arrow[dr, "q_{X_*}"] & \\
X_* & & X_*^\circ.
\end{tikzcd}
\]
\msp{page 140}
Looking at the restriction of the two morphisms of diagrams $(***)$ to the parts of the three diagrams of type $\Gamma^\circ$ over the tail $\{3,4\}$ of $\Gamma$ we find the following diagram in solid lines
\[
\begin{tikzcd}[column sep=3em, row sep=normal]
 & \int\Sigma(X_*) \arrow[dl, "W"'] \arrow[dr, "W"] \arrow[dd, "\varphi_{\Sigma(Q)}"'] & & \Sigma\tsX \arrow[dl, dotted, "W"'] \arrow[dr, dotted, "W"] \arrow[dd, dotted, "\Sigma\tilde{\varphi}_Q"] & \\
\sX \arrow[dd, "\varphi_Q"'] & & \tsX^\circ \arrow[dd, "\varphi_{Q^\circ}=(\tilde{\varphi}_Q)^\circ"'] & & \tsX \arrow[dd, dotted, "\tilde{\varphi}_Q"] \\
 & \Sigma(Y) \arrow[dl, "W"'] \arrow[dr, "W"] & & \Sigma(Y) \arrow[dl, dotted, "W"'] \arrow[dr, dotted, "W"] & \\
Y & & Y^\circ & & Y
\end{tikzcd}
\]
\printedpage{96}
% head: A. Grothendieck : Dérivateurs, Ch. XII, Édition des Universités de Montpellier II et Paris VII
% foot: Version 7 décembre 2002
which I have completed with a second panel in dotted lines. In this diagram of four adjoining squares, the horizontal sides of the squares represent arrows of $\Cat$ which are in $W$. We thus find: % [sic] qu j'ai -> que j'ai, missing letter
\[
\varphi_Q\in W \iff \varphi_{\Sigma Q}\in W \iff (\varphi_{Q^\circ}=\tilde{\varphi}_Q^\circ)\in W \iff \Sigma\tilde{\varphi}_Q\in W \iff \tilde{\varphi}_Q\in W.
\]
We also find that the five vertical arrows are isomorphic, as arrows of $\Hot_W$. To summarize:

\textbf{Proposition 14.} \textit{Let us consider a diagram (commutative square)}
\[
Q \quad = \qquad
\begin{tikzcd}[row sep=small, column sep=small]
 & X_0 \arrow[dl] \arrow[dr] & \\
X_1 \arrow[dr] & & X_2 \arrow[dl] \\
 & Y &
\end{tikzcd}
\]
\textit{in} $\Cat$\textit{, and the ``opposite'' diagram deduced by} $X\longmapsto X^\circ$
\[
Q^\circ \quad = \qquad
\begin{tikzcd}[row sep=small, column sep=small]
 & X_0^\circ \arrow[dl] \arrow[dr] & \\
X_1^\circ \arrow[dr] & & X_2^\circ \arrow[dl] \\
 & Y^\circ, &
\end{tikzcd}
\]
\msp{page 141}
\textit{and the corresponding arrows}
\[
\left\{
\begin{array}{rlclcl}
\varphi_Q : & \sX_Q & = &
\left(\displaystyle\int
\begin{tikzcd}[ampersand replacement=\&, row sep=small, column sep=small]
 \& X_1 \\
X_0 \arrow[ur] \arrow[dr] \& \\
 \& X_2
\end{tikzcd}
\right)
& \longrightarrow & Y \\[2ex]
\varphi_{Q^\circ} : & \sX_{Q^\circ} & = &
\left(\displaystyle\int
\begin{tikzcd}[ampersand replacement=\&, row sep=small, column sep=small]
 \& X_1^\circ \\
X_0^\circ \arrow[ur] \arrow[dr] \& \\
 \& X_2^\circ
\end{tikzcd}
\right)
& \longrightarrow & Y^\circ.
\end{array}
\right.
\]
\textit{Finally, let us also consider the canonical arrows $\tilde{\varphi}_Q$, $\tilde{\varphi}_{Q^\circ}$, relative to the proper integral $\int_{\mathrm{pr}}$} % [sic] relative -> relatives, agreement with « les flèches »
\printedpage{97}
% head: A. Grothendieck : Dérivateurs, Ch. XII, Édition des Universités de Montpellier II et Paris VII
% foot: Version 7 décembre 2002
\textit{(instead of the smooth integral $\int=\int_{\mathrm{lis}}$)}
\[
\left\{
\begin{array}{rlclcl}
\tilde{\varphi}_Q : & \tsX_Q & = &
\left(\displaystyle\int_{\mathrm{pr}}
\begin{tikzcd}[ampersand replacement=\&, row sep=small, column sep=small]
 \& X_1 \\
X_0 \arrow[ur] \arrow[dr] \& \\
 \& X_2
\end{tikzcd}
\right)
& \longrightarrow & Y \\[2ex]
\tilde{\varphi}_{Q^\circ} : & \tsX_{Q^\circ} & = &
\left(\displaystyle\int_{\mathrm{pr}}
\begin{tikzcd}[ampersand replacement=\&, row sep=small, column sep=small]
 \& X_1^\circ \\
X_0^\circ \arrow[ur] \arrow[dr] \& \\
 \& X_2^\circ
\end{tikzcd}
\right)
& \longrightarrow & Y^\circ.
\end{array}
\right.
\]
\textit{We then have the following:}

\begin{enumerate}[label={}, leftmargin=2.5em]\itshape
\item[\circled{\textit{a}}] $\varphi_{Q^\circ}\simeq(\tilde{\varphi}_Q)^\circ$, $\varphi_Q\simeq(\tilde{\varphi}_{Q^\circ})^\circ$ (from which it already follows that $\left\{\begin{array}{lcl}\varphi_Q\in W & \Longleftrightarrow & \tilde{\varphi}_{Q^\circ}\in W\\ \varphi_{Q^\circ}\in W & \Longleftrightarrow & \tilde{\varphi}_Q\in W\end{array}\right.$ (\footnote[90]{\upshape and that the arrows $\varphi_Q$ and $\tilde{\varphi}_{Q^\circ}$ on the one hand, $\varphi_{Q^\circ}$ and $\tilde{\varphi}_Q$ on the other hand, are isomorphic as arrows in $\Hot_W$.})).

\item[\circled{\textit{b}}] The relations
\[
\varphi_Q\in W, \qquad \varphi_{Q^\circ}\in W, \qquad \tilde{\varphi}_Q\in W, \qquad \tilde{\varphi}_{Q^\circ}\in W
\]
are equivalent, in other words: for the square to be $W$-homotopically cocartesian (smooth integral version), it is necessary and sufficient that the opposite square $Q^\circ$ be so, and these conditions are also equivalent to the corresponding conditions in terms of the proper integrals (instead of the smooth integrals).

\item[\circled{\textit{c}}] The four arrows $\varphi_Q$, $\varphi_{Q^\circ}$, $\tilde{\varphi}_Q$, $\tilde{\varphi}_{Q^\circ}$ of $\Cat$ as arrows in $\Hot_W$, are isomorphic in $\underline{\Fl}(\Hot_W)$. More precisely, as arrows of $\Cat$,
\msp{page 142}
they define isomorphic objects of $\HOT_W(\boldsymbol{\Delta}^1)$.

\item[\circled{\textit{d}}] More precisely still, the four diagrams in $\Cat$ of type $\Gamma^\circ$, where
\[
\Gamma = \;
\begin{tikzcd}[row sep=tiny, column sep=small]
 & 1 \arrow[dl, "\alpha_1"'] & & \\
0 & & 3 \arrow[ul, "i_1"'] \arrow[dl, "i_2"] & 4, \arrow[l, "\varphi"'] \\
 & 2 \arrow[ul, "\alpha_2"] & &
\end{tikzcd}
\]
$\alpha_1i_1\varphi=\alpha_2i_2\varphi$, into which the four arrows we have just considered fit, namely
\[
\widehat{Q} = 
\begin{tikzcd}[row sep=tiny, column sep=small]
 & X_1 \arrow[dr, "i_1"] & & \\
X_0 \arrow[ur, "\alpha_1"] \arrow[dr, "\alpha_2"'] \arrow[rr, dotted] & & \sX_Q \arrow[r, "\varphi_Q"] & Y, \\
 & X_2 \arrow[ur, "i_2"'] & &
\end{tikzcd}
\]
\end{enumerate}